\documentclass{amsart}
\usepackage{amsmath,amssymb,amsthm,hyperref}
\newtheorem{theorem}{Theorem}
\newtheorem{lemma}{Lemma}
\newtheorem{proposition}{Proposition}
\theoremstyle{definition}
\newtheorem{definition}{Definition}
\theoremstyle{remark}
\newtheorem{remark}{Remark}
\begin{document}

\title{An explicit recursive description of the generalized Cantor set}
\maketitle

\section{Setup and notation}

Let $(\alpha_n)_{n\ge1}$ satisfy $0<\alpha_n<1$ for all $n$ and the admissibility
condition
\begin{equation}\label{eq:adm}
1-\sum_{n=1}^{\infty}2^{n-1}\alpha_n\ \ge\ 0 .
\end{equation}
The set $C=\bigcap_{n\ge0}C_n$ is built by the middle--deletion procedure of the
problem statement. Throughout, $\{0,1\}^n$ denotes the set of binary strings of
length $n$ (with $\{0,1\}^0=\{()\}$, the empty string), and $\{0,1\}^{\mathbb N}$
denotes the set of infinite binary sequences $\varepsilon=(\varepsilon_k)_{k\ge1}$.
For $\varepsilon\in\{0,1\}^{\mathbb N}$ and $n\ge0$, $\varepsilon|n:=(\varepsilon_1,\dots,\varepsilon_n)$
is its length-$n$ prefix.

We first record the side length of the stage--$n$ intervals.

\begin{lemma}\label{lem:len}
For every $n\ge0$, each of the $2^{n}$ closed intervals making up $C_n$ has the
same length
\begin{equation}\label{eq:Ln}
L_n\ :=\ \frac{1}{2^{n}}\Bigl(1-\sum_{k=1}^{n}2^{k-1}\alpha_k\Bigr),
\end{equation}
with the convention that the empty sum is $0$, so $L_0=1$. The lengths obey the
recursion
\begin{equation}\label{eq:rec}
L_n=\frac{L_{n-1}-\alpha_n}{2}\qquad(n\ge1).
\end{equation}
Moreover $L_n>0$ for every finite $n$, so the construction is well defined at every
stage; the sequence $(L_n)_{n\ge0}$ is strictly decreasing, and
\begin{equation}\label{eq:Linf}
L_\infty\ :=\ \lim_{n\to\infty}L_n\ =\ 0 .
\end{equation}
\end{lemma}

\begin{proof}
\emph{Well-definedness and the recursion.}
We argue by induction, simultaneously proving \eqref{eq:Ln} and that the stage-$n$
deletion can actually be performed (each stage-$(n-1)$ interval is long enough to
contain an open middle of length $\alpha_n$). For $n=0$, $C_0=[0,1]$ is one
interval of length $L_0=1>0$, matching \eqref{eq:Ln}. Assume $C_{n-1}$ consists of
$2^{n-1}$ closed intervals each of length $L_{n-1}$, where \eqref{eq:Ln} holds for
$n-1$.

We first check $L_{n-1}>\alpha_n$, which is exactly what is needed to remove an open
middle of (absolute) length $\alpha_n$ from an interval of length $L_{n-1}$ and be
left with two nondegenerate closed pieces. By the admissibility condition
\eqref{eq:adm} the full series satisfies $\sum_{k=1}^{\infty}2^{k-1}\alpha_k\le1$;
since all terms are positive, the partial sum through index $n$ is \emph{strictly}
less than $1$:
\[
\sum_{k=1}^{n}2^{k-1}\alpha_k
\ <\ \sum_{k=1}^{n}2^{k-1}\alpha_k+\sum_{k=n+1}^{\infty}2^{k-1}\alpha_k
\ =\ \sum_{k=1}^{\infty}2^{k-1}\alpha_k\ \le\ 1,
\]
where the strict inequality uses $2^{n}\alpha_{n+1}>0$ (the tail is a sum of
positive terms, hence positive). Consequently, by \eqref{eq:Ln} for $n-1$,
\[
L_{n-1}-\alpha_n
=\tfrac1{2^{n-1}}\Bigl(1-\sum_{k=1}^{n-1}2^{k-1}\alpha_k\Bigr)-\alpha_n
=\tfrac1{2^{n-1}}\Bigl(1-\sum_{k=1}^{n}2^{k-1}\alpha_k\Bigr)\ >\ 0,
\]
i.e.\ $L_{n-1}>\alpha_n$. Thus at stage $n$ each stage-$(n-1)$ interval has its open
middle subinterval of (absolute) length $\alpha_n$ removed; this leaves two closed
intervals, one on each side, of equal length $L_n=(L_{n-1}-\alpha_n)/2>0$, which is
\eqref{eq:rec}. Hence $C_n$ is a union of $2^{n}$ closed intervals of common
positive length $L_n$. Solving \eqref{eq:rec} with $L_0=1$: if \eqref{eq:Ln} holds
for $n-1$ then
\[
L_n=\tfrac12 L_{n-1}-\tfrac12\alpha_n
=\tfrac{1}{2^{n}}\Bigl(1-\sum_{k=1}^{n-1}2^{k-1}\alpha_k\Bigr)-\tfrac{\alpha_n}{2}
=\tfrac{1}{2^{n}}\Bigl(1-\sum_{k=1}^{n}2^{k-1}\alpha_k\Bigr),
\]
proving \eqref{eq:Ln} for $n$ and completing the induction.

\smallskip
\emph{Monotonicity and the limit.}
By \eqref{eq:rec}, $L_{n-1}-L_n=\tfrac12(L_{n-1}+\alpha_n)>0$ because $L_{n-1}>0$ and
$\alpha_n>0$, so $(L_n)$ is strictly decreasing; bounded below by $0$, it converges.
To identify the limit, write $S_n:=\sum_{k=1}^{n}2^{k-1}\alpha_k$ and
$S:=\sum_{k=1}^{\infty}2^{k-1}\alpha_k\le1$ by \eqref{eq:adm}. Then by \eqref{eq:Ln}
\[
0\ <\ L_n\ =\ 2^{-n}\,(1-S_n)\ \le\ 2^{-n}\cdot 1\ =\ 2^{-n},
\]
since $0\le 1-S_n\le 1$ (the factor $1-S_n$ is bounded, indeed $1-S_n\downarrow 1-S$
is a fixed finite number). As $2^{-n}\to0$, the squeeze gives $L_n\to0$, i.e.\
$L_\infty=0$, proving \eqref{eq:Linf}.
\end{proof}

\begin{remark}[The limit length is $0$, the total measure is $1-S$]
It is essential to distinguish two quantities. The \emph{per-interval} length is
$L_n=2^{-n}(1-S_n)\to0$. The \emph{total} length of $C_n$ is
\[
|C_n|\ =\ 2^{n}L_n\ =\ 1-S_n\ \longrightarrow\ 1-S\ =\ 1-\sum_{k=1}^{\infty}2^{k-1}\alpha_k .
\]
The factor $2^{-n}$ counting the per-interval scale must not be cancelled against
the $2^{n}$ counting the number of intervals: $L_\infty=0$ always, whereas the
Lebesgue measure $|C|=1-S$ may be positive. The measure of $C$ is computed
separately in Theorem~\ref{thm:meas} below by downward continuity of Lebesgue
measure on the nested compacts $C_n$.
\end{remark}

It is convenient to introduce the \emph{step sizes}
\begin{equation}\label{eq:delta}
\delta_k\ :=\ L_{k-1}-L_k\ =\ \tfrac12\bigl(L_{k-1}+\alpha_k\bigr)\ >\ 0
\qquad(k\ge1),
\end{equation}
where the middle equality is \eqref{eq:rec} and positivity holds because
$L_{k-1}>0$ and $\alpha_k>0$. Telescoping \eqref{eq:delta} and using $L_\infty=0$
gives, for every $n\ge0$,
\begin{equation}\label{eq:tail}
\sum_{k=n+1}^{\infty}\delta_k\ =\ L_n-L_\infty\ =\ L_n,
\qquad\text{in particular}\qquad
\sum_{k=1}^{\infty}\delta_k=L_0=1<\infty .
\end{equation}

\section{Labelling the intervals; the address map}

We label the $2^n$ intervals of $C_n$ by binary strings of length $n$ compatibly
with the parent--child relation: when the interval labelled
$\varepsilon=(\varepsilon_1,\dots,\varepsilon_n)$ is split at stage $n+1$, its
\emph{left} child gets the label $\varepsilon0$ and its \emph{right} child gets
$\varepsilon1$. Write $I_\varepsilon=[x_\varepsilon,\,x_\varepsilon+L_n]$ for the
interval labelled $\varepsilon\in\{0,1\}^n$.

\begin{proposition}\label{prop:endpoint}
For every $n\ge0$ and every $\varepsilon=(\varepsilon_1,\dots,\varepsilon_n)\in\{0,1\}^n$,
\begin{equation}\label{eq:xeps}
x_\varepsilon\ =\ \sum_{k=1}^{n}\varepsilon_k\,\delta_k ,
\qquad\text{so}\qquad
I_\varepsilon=\Bigl[\sum_{k=1}^{n}\varepsilon_k\delta_k,\ \
\sum_{k=1}^{n}\varepsilon_k\delta_k+L_n\Bigr].
\end{equation}
Moreover the two children of $I_\varepsilon$ are
\begin{equation}\label{eq:children}
I_{\varepsilon0}=[\,x_\varepsilon,\ x_\varepsilon+L_{n+1}\,],
\qquad
I_{\varepsilon1}=[\,x_\varepsilon+\delta_{n+1},\ x_\varepsilon+\delta_{n+1}+L_{n+1}\,],
\end{equation}
and they are separated by the deleted open gap $(x_\varepsilon+L_{n+1},\,x_\varepsilon+\delta_{n+1})$ of length $\alpha_{n+1}>0$.
\end{proposition}

\begin{proof}
Induction on $n$. For $n=0$ the only string is empty and $x_{()}=0$, the left
endpoint of $C_0=[0,1]$. Assume \eqref{eq:xeps} for length $n$, and fix
$\varepsilon\in\{0,1\}^n$ with $I_\varepsilon=[x_\varepsilon,x_\varepsilon+L_{n}]$.
Deleting the open middle of length $\alpha_{n+1}$ leaves two closed intervals each
of length $L_{n+1}$: the left one occupies $[x_\varepsilon,x_\varepsilon+L_{n+1}]$,
then a gap of length $\alpha_{n+1}$, then the right one of length $L_{n+1}$; since
$L_{n+1}+\alpha_{n+1}+L_{n+1}=L_n$ by \eqref{eq:rec}, the right one is
$[x_\varepsilon+L_{n+1}+\alpha_{n+1},\,x_\varepsilon+L_n]$. Using
$\delta_{n+1}=L_n-L_{n+1}$ from \eqref{eq:delta} together with
$L_n=2L_{n+1}+\alpha_{n+1}$ from \eqref{eq:rec},
\[
L_{n+1}+\alpha_{n+1}=L_{n+1}+(L_n-2L_{n+1})=L_n-L_{n+1}=\delta_{n+1},
\]
which gives \eqref{eq:children} and the stated gap, of length
$\delta_{n+1}-L_{n+1}=\alpha_{n+1}>0$. Hence
\[
x_{\varepsilon0}=x_\varepsilon=\sum_{k=1}^{n}\varepsilon_k\delta_k,
\qquad
x_{\varepsilon1}=x_\varepsilon+\delta_{n+1}=\sum_{k=1}^{n+1}\varepsilon_k\delta_k
\ \ (\text{coordinate }n{+}1\text{ equal to }1),
\]
so $x_{\varepsilon'}=\sum_{k=1}^{n+1}\varepsilon'_k\delta_k$ for every length-$(n+1)$
string $\varepsilon'$, completing the induction.
\end{proof}

For $\varepsilon\in\{0,1\}^{\mathbb N}$ define the \emph{address value}
\begin{equation}\label{eq:phi}
\varphi(\varepsilon)\ :=\ \sum_{k=1}^{\infty}\varepsilon_k\,\delta_k ,
\end{equation}
which converges since $0\le\varepsilon_k\delta_k\le\delta_k$ and
$\sum_k\delta_k<\infty$ by \eqref{eq:tail}.

\section{Main theorem: the explicit description of $C$}

\begin{theorem}\label{thm:main}
Set $L_0=1$, $L_k=\tfrac12(L_{k-1}-\alpha_k)$, and
$\delta_k=L_{k-1}-L_k=\tfrac12(L_{k-1}+\alpha_k)$, so that explicitly
$L_n=2^{-n}\bigl(1-\sum_{k=1}^{n}2^{k-1}\alpha_k\bigr)$. Then the constructed set
admits the explicit closed form
\begin{equation}\label{eq:main}
\boxed{\,C\ =\ \Bigl\{\ \sum_{k=1}^{\infty}\varepsilon_k\,\delta_k\ :\
(\varepsilon_k)_{k\ge1}\in\{0,1\}^{\mathbb N}\ \Bigr\}
\ =\ \varphi\bigl(\{0,1\}^{\mathbb N}\bigr),\,}
\end{equation}
and the address map $\varphi:\{0,1\}^{\mathbb N}\to C$ of \eqref{eq:phi} is a
bijection. This holds in all admissible cases \eqref{eq:adm}, regardless of whether
$\sum_{k}2^{k-1}\alpha_k$ equals $1$ or is strictly less than $1$.
\end{theorem}

\begin{proof}
Recall $C_n=\bigcup_{\varepsilon\in\{0,1\}^n}I_\varepsilon$ and
$C=\bigcap_{n\ge0}C_n$, with $I_\varepsilon$ as in \eqref{eq:xeps}.

\smallskip
\emph{Step 0 (nesting and disjointness of the dyadic intervals).}
By \eqref{eq:children}, each child interval $I_{\varepsilon b}$ ($b\in\{0,1\}$) is
contained in its parent $I_\varepsilon$; hence for every
$\varepsilon\in\{0,1\}^{\mathbb N}$ the prefix intervals are nested,
$I_{\varepsilon|0}\supseteq I_{\varepsilon|1}\supseteq\cdots$. Also by
\eqref{eq:children} the two children of any interval are separated by a gap of
length $\alpha_{n+1}>0$, hence are disjoint. Consequently any two distinct strings
$\sigma\ne\sigma'$ of the same length $n$ have $I_\sigma\cap I_{\sigma'}=\varnothing$:
let $m\le n$ be the first index where they differ, write $\rho=\sigma|(m-1)=\sigma'|(m-1)$
for their common length-$(m-1)$ prefix, and assume $\sigma_m=0,\sigma'_m=1$ (the
other case is symmetric). By the nesting just established applied to the suffixes,
$I_\sigma\subseteq I_{\sigma|m}=I_{\rho 0}$ and $I_{\sigma'}\subseteq
I_{\sigma'|m}=I_{\rho 1}$; since $I_{\rho 0}$ and $I_{\rho 1}$ are the two
(disjoint) children of $I_\rho$, the inclusions force $I_\sigma\cap
I_{\sigma'}=\varnothing$.

\smallskip
\emph{Step 1 (each branch shrinks to a single point).}
Fix $\varepsilon\in\{0,1\}^{\mathbb N}$. The intervals $I_{\varepsilon|n}$ are
nested and closed, with lengths $L_n\to0$ by \eqref{eq:Linf}. Their left endpoints
$x_{\varepsilon|n}=\sum_{k=1}^{n}\varepsilon_k\delta_k$ increase to
$\varphi(\varepsilon)$, while their right endpoints
$x_{\varepsilon|n}+L_n\to\varphi(\varepsilon)+L_\infty=\varphi(\varepsilon)$. By the
nested interval theorem (or Cantor's intersection theorem) the intersection is a
single point:
\begin{equation}\label{eq:cell}
\bigcap_{n\ge0}I_{\varepsilon|n}=\{\varphi(\varepsilon)\}.
\end{equation}
(That the right endpoints do not increase along the branch follows from
\eqref{eq:children}: passing from $I_{\varepsilon|n}$ to $I_{\varepsilon|(n+1)}$
either keeps the left endpoint with new right endpoint $x_{\varepsilon|n}+L_{n+1}\le
x_{\varepsilon|n}+L_n$ when $\varepsilon_{n+1}=0$, or moves to right endpoint
$x_{\varepsilon|n}+\delta_{n+1}+L_{n+1}=x_{\varepsilon|n}+L_n$ when
$\varepsilon_{n+1}=1$; in either case the nesting is confirmed.)

\smallskip
\emph{Step 2 ($C=\varphi(\{0,1\}^{\mathbb N})$).}
($\supseteq$) For any $\varepsilon$ and any $n$,
$\varphi(\varepsilon)\in I_{\varepsilon|n}\subseteq C_n$; hence
$\varphi(\varepsilon)\in\bigcap_n C_n=C$. Thus
$\varphi(\{0,1\}^{\mathbb N})\subseteq C$.

($\subseteq$) Let $t\in C$. For each $n$, $t\in C_n$, so $t\in I_{\sigma^{(n)}}$
for a string $\sigma^{(n)}\in\{0,1\}^n$ that is \emph{unique} by the disjointness
in Step 0. If $\tau$ denotes the length-$n$ prefix of $\sigma^{(n+1)}$, then
$t\in I_{\sigma^{(n+1)}}\subseteq I_\tau$ (nesting), so uniqueness at level $n$
forces $\tau=\sigma^{(n)}$; thus the $\sigma^{(n)}$ are nested prefixes and define
a single $\varepsilon\in\{0,1\}^{\mathbb N}$ with $\varepsilon|n=\sigma^{(n)}$ for
all $n$. Hence $t\in\bigcap_n I_{\varepsilon|n}=\{\varphi(\varepsilon)\}$ by
\eqref{eq:cell}, i.e.\ $t=\varphi(\varepsilon)$. Therefore $C\subseteq
\varphi(\{0,1\}^{\mathbb N})$, and \eqref{eq:main} is established.

\smallskip
\emph{Step 3 ($\varphi$ is a bijection onto $C$).}
Surjectivity onto $C$ is \eqref{eq:main}. For injectivity, let
$\varepsilon\ne\varepsilon'$ and let $m\ge1$ be the first index where they differ;
say $\varepsilon_m=0,\ \varepsilon'_m=1$ (the other case is symmetric). Using
$\varepsilon'_k-\varepsilon_k\ge-1$ together with \eqref{eq:tail},
\[
\varphi(\varepsilon')-\varphi(\varepsilon)
=\delta_m+\sum_{k=m+1}^{\infty}(\varepsilon'_k-\varepsilon_k)\delta_k
\ \ge\ \delta_m-\sum_{k=m+1}^{\infty}\delta_k
=\delta_m-L_m .
\]
Using $\delta_m=L_{m-1}-L_m$ from \eqref{eq:delta} and $L_{m-1}=2L_m+\alpha_m$ from
\eqref{eq:rec},
\[
\delta_m-L_m=(L_{m-1}-L_m)-L_m=L_{m-1}-2L_m=\alpha_m\ >\ 0,
\]
so $\varphi(\varepsilon')>\varphi(\varepsilon)$. In particular
$\varphi(\varepsilon)\ne\varphi(\varepsilon')$, so $\varphi$ is injective, hence a
bijection $\{0,1\}^{\mathbb N}\to C$. (The estimate moreover shows $\varphi$ is
strictly order preserving with respect to the lexicographic order on
$\{0,1\}^{\mathbb N}$ and the usual order on $\mathbb R$.)
\end{proof}

\section{Topological properties and Lebesgue measure}

\begin{proposition}\label{prop:topo}
$C$ is a nonempty compact, perfect, totally disconnected, nowhere dense subset of
$[0,1]$.
\end{proposition}

\begin{proof}
\emph{Nonempty and compact.} $C=\bigcap_n C_n$ is an intersection of nonempty
nested compact sets, hence nonempty and compact.

\smallskip
\emph{Totally disconnected.} Let $s=\varphi(\varepsilon)$ and
$t=\varphi(\varepsilon')$ be two distinct points of $C$ with $s<t$, and let $m$ be
the first index where the addresses differ. By the order preservation in Step~3 we
must have $\varepsilon_m=0,\ \varepsilon'_m=1$. Both points lie in the common
ancestor interval $I_{\varepsilon|(m-1)}$, but
$s=\varphi(\varepsilon)\in I_{(\varepsilon|(m-1))0}=I_{\varepsilon|m}$ and
$t=\varphi(\varepsilon')\in I_{(\varepsilon|(m-1))1}$, the two children of
$I_{\varepsilon|(m-1)}$. These children are separated by the deleted open gap
$G:=(x_{\varepsilon|(m-1)}+L_m,\ x_{\varepsilon|(m-1)}+\delta_m)$ of length
$\alpha_m>0$ (Proposition~\ref{prop:endpoint}), and $G\cap C=\varnothing$ because
$G$ is disjoint from $C_m\supseteq C$. Pick any $c\in G$; then $s<c<t$ and
$c\notin C$, so $s$ and $t$ lie in different connected components of $C$. As $s,t$
were arbitrary distinct points, no connected subset of $C$ contains two points;
i.e.\ $C$ is totally disconnected.

\smallskip
\emph{Perfect (no isolated points).} Fix $\varepsilon\in\{0,1\}^{\mathbb N}$ and
$s=\varphi(\varepsilon)$. For each $n$, let $\varepsilon^{(n)}$ be $\varepsilon$
with its $n$-th coordinate flipped; then $\varepsilon^{(n)}\ne\varepsilon$ but
$\varepsilon^{(n)}$ agrees with $\varepsilon$ on all coordinates $>n$ except (they
also agree below $n$). Their first difference is at index $n$, so by the bound in
Step~3 (applied to whichever of $s,\varphi(\varepsilon^{(n)})$ is larger),
\[
|\varphi(\varepsilon^{(n)})-\varphi(\varepsilon)|
=\Bigl|\sum_{k=1}^{\infty}\varepsilon^{(n)}_k\delta_k-\sum_{k=1}^{\infty}\varepsilon_k\delta_k\Bigr|
=\delta_n\ \le\ L_{n-1}\ \le\ 2^{-(n-1)}\ \longrightarrow\ 0,
\]
where we used $\delta_n=L_{n-1}-L_n\le L_{n-1}$ and $L_{n-1}\le 2^{-(n-1)}$ from
Lemma~\ref{lem:len}. (Indeed $|\varphi(\varepsilon^{(n)})-\varphi(\varepsilon)|
=\delta_n$ exactly, since the two addresses differ only in coordinate $n$.) Thus
$\varphi(\varepsilon^{(n)})\in C\setminus\{s\}$ converges to $s$, so $s$ is not
isolated. As $C$ is also closed, it is perfect.

\smallskip
\emph{Nowhere dense.} $C$ is closed, so it suffices to show it has empty interior.
If $C$ contained an interval $J$ of positive length $\ell$, then $J\subseteq C_n$
for all $n$, forcing $\ell\le L_n\to0$, a contradiction. (Total disconnectedness
already gives empty interior; we record the direct argument for completeness.)
\end{proof}

\begin{theorem}\label{thm:meas}
The Lebesgue measure of $C$ is
\begin{equation}\label{eq:measure}
|C|\ =\ 1-\sum_{k=1}^{\infty}2^{k-1}\alpha_k\ =\ 1-S .
\end{equation}
In particular, when the admissibility inequality \eqref{eq:adm} is strict
($S<1$), $C$ is a \emph{fat} Cantor set: a totally disconnected, perfect, nowhere
dense set of \emph{positive} Lebesgue measure $1-S$ (a Smith--Volterra--Cantor type
set). When $S=1$, $C$ has measure $0$.
\end{theorem}

\begin{proof}
Each $C_n$ is a finite disjoint union of $2^n$ closed intervals of length $L_n$, so
\[
|C_n|\ =\ 2^{n}L_n\ =\ 2^{n}\cdot 2^{-n}(1-S_n)\ =\ 1-S_n,
\qquad S_n=\sum_{k=1}^{n}2^{k-1}\alpha_k .
\]
The sets $C_n$ are nested decreasing compacts with $|C_0|=1<\infty$, and
$C=\bigcap_n C_n$. By downward continuity of Lebesgue measure for a decreasing
sequence of measurable sets of finite measure,
\[
|C|\ =\ \lim_{n\to\infty}|C_n|\ =\ \lim_{n\to\infty}(1-S_n)\ =\ 1-S
\ =\ 1-\sum_{k=1}^{\infty}2^{k-1}\alpha_k,
\]
which is \eqref{eq:measure}. The dichotomy ($1-S>0$ vs.\ $1-S=0$) follows directly
from \eqref{eq:adm}.
\end{proof}

\begin{remark}[The per-interval length and the measure are different quantities]
Theorem~\ref{thm:meas} computes the \emph{total} measure $1-S$ as
$\lim_n 2^{n}L_n$, while the \emph{per-interval} length is $L_n\to L_\infty=0$.
These must not be conflated: the explicit description \eqref{eq:main} writes $C$ as
the image of the address map $\varphi$, never as a union of nondegenerate
intervals. (An uncountable disjoint union of intervals of a common positive length
cannot fit in $[0,1]$, and would contradict total disconnectedness; no such
structure occurs.) Even when $|C|=1-S>0$, the set $C$ is the totally disconnected
image $\varphi(\{0,1\}^{\mathbb N})$ of the Cantor space.
\end{remark}

\begin{remark}[Classical middle--thirds set]
Taking $\alpha_n=3^{-n}$ gives $\sum_{n\ge1}2^{n-1}\alpha_n=\tfrac13\sum_{n\ge1}(2/3)^{n-1}=1$,
so $S=1$ and $|C|=0$, and one computes $\delta_k=L_{k-1}-L_k=2\cdot3^{-k}$. Then
\eqref{eq:main} reads
$C=\bigl\{\sum_{k\ge1}2\varepsilon_k 3^{-k}:\varepsilon_k\in\{0,1\}\bigr\}$, the
standard base-$3$ description of the Cantor middle--thirds set.
\end{remark}

\end{document}
